% ============================================================================
% Fichier cible : ARIEL_LOGMINER_MEMOIRE_FINAL/chapters/chapitre2_etat_art.tex
% Passage à remplacer : de « \section{Lacunes identifiées dans la littérature} »
%   (inclus) à « \section{Synthèse du chapitre} » (exclu).
% Agent : Agent 3 — État de l'art et positionnement, chantier 2026-11-05.
% Sources : état de l'art du chapitre ; cadres intégrés et schémas ajoutés par
%   l'Agent 3 ; résultats du corps ; C12 et fichier Agent 2
%   chapitre2_etat_art.tex__tableau_1.3_positionnement.tex.
% ============================================================================

\section{Tensions de la littérature et espace expérimental}

L'état de l'art n'établit pas l'absence d'une chaîne intégrée. Wazuh, Elastic Security et
OpenSearch Security Analytics couvrent déjà ingestion, règles, alertes, investigation ou
corrélation à des degrés différents
\cite{wazuh2026architecture,elastic2026security,opensearch2026securityanalytics}.
LogAI et Loglizer fournissent par ailleurs des cadres pour comparer des traitements de
journaux \cite{cheng2023logai,he2016experience}. La lacune initialement formulée comme un
« espace entre prototype algorithmique et plateforme opérationnelle » ne suffit donc pas
à établir une nouveauté.

La tension pertinente est méthodologique. Les travaux algorithmiques isolent souvent le
détecteur afin d'en mesurer la performance, alors que les plateformes intégrées offrent
de nombreuses fonctions mais ne constituent pas nécessairement un protocole de recherche
contrôlant unité d'analyse, partition, baseline, routage et coût d'orchestration. Cette
tension ouvre un espace pour un banc d'expérimentation instrumenté, à condition de ne pas
confondre intégration et contribution scientifique.

Quatre exigences structurent alors une comparaison défendable :
\begin{itemize}
    \item mesurer le mapping vers un schéma reconnu, plutôt que déclarer une enveloppe
          locale « normalisée » par construction ;
    \item comparer le routage et Contract Net à une table statique ou à des files
          spécialisées, et non seulement à un monolithe ;
    \item évaluer les modèles contre des règles triviales sur la même unité et le même
          test, avec des partitions qui exposent le transfert ;
    \item distinguer les propriétés fournies par le broker, par le magasin d'idempotence,
          par le code des agents et par le modèle prédictif.
\end{itemize}

Ariel Logminer satisfait partiellement cette ambition : les artefacts relient plusieurs
affirmations à des protocoles, et les résultats négatifs révèlent les raccourcis de
datasets. En revanche, aucune comparaison expérimentale n'est conduite contre les cadres
intégrés, aucun incident réel ne combine journaux système et flux réseau, et le bénéfice
propre de la négociation n'est pas établi dans le manuscrit actuel. Les observations sur
la dépendance aux partitions répliquent des avertissements déjà publiés
\cite{sommer2010outside,arp2022dos,le2022logbased,engelen2021troubleshooting,liu2022error} ;
elles ne constituent pas une découverte conceptuelle.

\section{Positionnement scientifique d'Ariel Logminer}

\logminer{} est positionné comme un \textbf{prototype d'intégration instrumenté pour
l'étude critique de chaînes légères d'analyse de journaux}. Il n'introduit ni algorithme
de détection, ni schéma de normalisation, ni protocole d'allocation nouveau. Son apport
potentiel réside dans la mise en relation, au sein d'un même prototype, de choix
d'ingénierie et de protocoles dont les effets peuvent être séparés : unité événement ou
bloc, partition aléatoire ou temporelle, règle triviale ou modèle, routage, orchestration,
idempotence et corrélation.

Ce positionnement est une contribution d'ingénierie expérimentale, dont la valeur
scientifique dépend des comparateurs réellement exécutés. Les cadres intégrés établissent
que l'assemblage fonctionnel n'est pas nouveau. Les expériences du mémoire établissent
pour leur part plusieurs limites : forte dépendance aux jeux et aux partitions, coût de
l'orchestration, distribution déséquilibrée, normalisation partielle et corrélation
validée seulement en circuit fermé. Le tableau~\ref{tab:positionnement-logminer} sépare
donc strictement les propriétés observées des revendications non démontrées.

\begin{table}[H]
\centering
\caption{Positionnement scientifique d'Ariel Logminer}
\label{tab:positionnement-logminer}
\small
\begin{tabularx}{\textwidth}{p{3.15cm}X X}
\toprule
\textbf{Dimension} &
\textbf{Ce que le prototype établit ou permet d'étudier} &
\textbf{Ce qu'il ne revendique pas} \\
\midrule
Hétérogénéité &
Lecture d'un ensemble documenté de formats et projection partielle dans un schéma commun ; chaque source est évaluée séparément. &
Compatibilité universelle avec tout format de journal ; analyse conjointe d'un même incident à partir de journaux système et de flux réseau. \\
Normalisation &
Contrat local extensible et mesure du remplissage de sept champs communs. &
Conformité ECS ou OCSF ; conservation sémantique complète de chaque source. \\
Détection &
Comparaison de modèles et de protocoles sur plusieurs jeux publics. &
Supériorité générale sur les règles triviales ou transfert à un environnement opérationnel. \\
Routage &
Orientation de fichiers connus et rejet open-set sous un seuil gelé. &
Gain prédictif systématique ; avantage établi sur une table de routage statique. \\
Agents &
Réalisation d'échanges inspirés de Contract Net, état persistant et décisions locales bornées. &
Autonomie cognitive, apprentissage de la politique ou supériorité de la négociation. \\
Distribution &
Fonctionnement multiprocessus et campagnes multi-VM de laboratoire. &
Équilibrage démontré, scalabilité industrielle, multi-site ou haute disponibilité complète. \\
Corrélation &
Regroupement temporel et contextuel évalué sur des corpus synthétiques contrôlés. &
Reconstruction causale ou validation sur incidents SOC réels multi-sources. \\
Mémoire &
Persistance de l'historique et support d'une procédure de mise à jour contrôlée. &
Gain mesuré de la mémoire sur l'attribution des tâches ; apprentissage continu autonome en production. \\
Priorisation &
Mécanisme prévu par la conception. &
Amélioration mesurée de l'ordre de traitement des anomalies. \\
Gouvernance &
Registre et chemins de comparaison ou de promotion testables. &
Cycle MLOps validé de bout en bout, surveillance de dérive et retour arrière éprouvé. \\
Dashboard &
Restitution fonctionnelle des informations et support de validation. &
Amélioration du travail de l'analyste démontrée par une étude utilisateur. \\
\bottomrule
\end{tabularx}
\end{table}

